Use the deep orbit of each half as a reference for all its earlier
turns. Lu compares secants to a common terminal section by bounding
the variation of its derivative \cite[Lemma~3.4]{LuDegree}.
On a regular descent we can instead integrate the exact difference of
the two graph equations against the same reference leg. The remainder
in the unscaled logarithm is bounded uniformly in the two turns.

\begin{lemma}[Upper deficits on a regular descent]\label{lem:deficit-ratio}
Let $\eps\downarrow0$. For each scaled profile let $T_b$ be a later turn
and $I=[T_1,T_2]\subset(0,T_b)$ a compact interval of earlier
turns. All these outgoing legs have transverse positive upper hits.
Write $\mathfrak d(T)=X(T_b)-X(T)$. Suppose $|\phi'|\le M$ on $[0,T_b]$,
and on $[T_1,T_b]$ suppose $\nu=-\phi'$ satisfies
$0<\nu_*\le\nu\le\nu^*$ and $|\nu'|\le M_1$. Assume
$T_b-T_2\ge b>0$, $|I|\le L$, and assume the upper hits of the
reference and test legs lie in a compact subinterval of $(0,\infty)$.
For a family of profiles take all these bounds common to the family.
Then, uniformly for $T,T_0\in I$ and all sufficiently small $\eps$,
\begin{equation}\label{eq:deficit-ratio}
 \log\frac{\mathfrak d(T)}{\mathfrak d(T_0)}=-\frac{1}{\eps}\int_{T_0}^{T}
 (\phi')^2\dd u+O(1).
\end{equation}
\end{lemma}
\begin{proof}
Let $w_*$ be the outgoing leg of $T_b$, $\zeta=w_*-\phi$ and
$d_T=w_*-w_T>0$ up to the earlier turn. In backward distance
$v=T_b-u$ put $Z(v)=\zeta(T_b-v)$. Then $Z(0)=0$ and for $v>0$,
\[
 (Z^2)_v=2\eps+2\phi'(T_b-v)Z .
\]
The right side extends continuously to $v=0$. Integrating first from
a positive lower endpoint and taking its limit gives this identity
in integral form from the turn. Put $c=1/[2(M+1)]$. While $Z\le c\eps$
the squared gap grows at speed at least $\eps$. It reaches $c^2\eps^2$
by $v=c^2\eps$ and cannot recross this level downwards. For $\eps$ so
small that $c^2\eps<b/2$, we have $\zeta(u)\ge c\eps$ from every
test turn back to the axis.

On the descent the same squared equation at $Z=\eps/\nu_*$ gives an
upper barrier starting from zero. Write $N(v)=\nu(T_b-v)$ there.
The error $e=Z-\eps/N$ obeys $e_v=-(N/Z)e-(\eps/N)_v$, with
$N/Z\ge\nu_*^2/\eps$ and $|(\eps/N)_v|\le M_1\eps/\nu_*^2$.
At $v=b/2$, $|e|\le C\eps$. Hence for $T_b-v\in I$,
$|e|\le C\eps e^{-\nu_*^2 b/(2\eps)}+C\eps^2=O(\eps^2)$.

The exact difference equation yields an integrable identity even at
the earlier turn,
\[
 (\log d_T)'=\frac{\eps}{\zeta(\zeta-d_T)}
             =\frac{\eps}{\zeta^2}+\frac{d_T'}{\zeta},\qquad d_T(T)=\zeta(T).
\]
Indeed $d_T'>0$, so its contribution at the turn satisfies
\[
 0\le E_T:=\int_0^T\frac{d_T'}\zeta\,\dd u
 \le\frac{\zeta(T)-d_T(0)}{\inf_{[0,T]}\zeta}=O(1).
\]
The numerator is at most $\zeta(T)=O(\eps)$ and the denominator is at
least $c\eps$. Integrating up to the turn by a one-sided limit gives
\[
 \log d_T(0)=\log \zeta(T)-\int_0^T\frac{\eps}{\zeta^2}\,\dd u-E_T .
\]
The same reference leg $w_*$ occurs for $T_0$, hence subtraction
cancels the integral on the common interval exactly. Also
$\mathfrak d(T)=d_T(0)[w_*(0)+w_T(0)]/2$. The upper-hit bounds and the bounds
on $\zeta(T)/\eps$ give bounded logarithms for both endpoint ratios. Thus
\[
 \log\frac{\mathfrak d(T)}{\mathfrak d(T_0)}
 =-\int_{T_0}^T\frac{\eps}{\zeta^2}\,\dd u+O(1).
\]
On $I$ the expansion for $\zeta$ gives
$\eps/\zeta^2=\nu^2/\eps+O(1)$. The integration length is at most
$L$, which gives the uniform remainder in \eqref{eq:deficit-ratio}.
\end{proof}
